% Part: incompleteness
% Chapter: theories-computability
% Section: extensions-of-q-not-decidable

\documentclass[../../../include/open-logic-section]{subfiles}

\begin{document}

\olfileid[pt]{inc}{tcp}{cqn}

\olsection{Extensões consistentes de $\Th{Q}$ são indecidíveis}

\begin{explain}
Lembre-se de que uma teoria é \emph{consistente} se não prova tanto
$!A$ quanto $\lnot !A$ para nenhuma fórmula~$!A$. Como qualquer coisa
decorre de uma contradição, uma teoria inconsistente é trivial: toda
sentença é demonstrável. É claro que, se uma teoria é
$\omega$-consistente, então é consistente. Mas a consistência é uma
exigência mais fraca (O.s.\@, há teorias consistentes que não são
$\omega$-consistentes). Podemos enfraquecer a hipótese de
\olref[oqn]{thm:oconsis-q} para a simples consistência e obter um
teorema mais forte.
\end{explain}

\begin{lem}
Não existe uma ``relação computável universal''. Isto é, não existe
relação computável binária $R(x,y)$ com a seguinte propriedade: sempre
que $S(y)$ é uma relação computável unária, existe algum $k$ tal que,
para todo $y$, $S(y)$ é verdadeira se, e somente se, $R(k,y)$ é
verdadeira.
\end{lem}

\begin{proof}
Suponha que $R(x,y)$ seja uma relação computável universal. Seja $S(y)$
a relação $\lnot R(y,y)$. Como $S(y)$ é computável, para algum $k$,
$S(y)$ é equivalente a $R(k,y)$. Mas então $S(k)$ é equivalente tanto
a $R(k,k)$ quanto a $\lnot R(k,k)$, o que é uma contradição.
\end{proof}


\begin{thm}
Seja $\Th{T}$ uma teoria consistente qualquer que inclua $\Th{Q}$.
Então $\Th{T}$ não é decidível.
\end{thm}


\begin{proof}
Suponha que $\Th{T}$ seja uma extensão consistente e decidível de
$\Th{Q}$. Obteremos uma contradição usando $\Th{T}$ para definir uma
relação computável universal.

Faça $R(x,y)$ valer se, e somente se,
\begin{quote}
$x$ codifica uma fórmula $!D(u)$ e $\Th{T}$ prova $!D(\num y)$.
\end{quote}
Como estamos supondo que $\Th{T}$ é decidível, $R$ é computável.
Mostremos que $R$ é universal. Se $S(y)$ é uma relação computável
qualquer, então é representável em $\Th{Q}$ (e, portanto, em $\Th{T}$)
por uma fórmula $!D_S(u)$. Então, para todo $n$, temos
\begin{eqnarray*}
S(\num n) & \lif & T \vdash !D_S(\num n) \\
& \lif & R(\Gn{!D_S(u)}, n)
\end{eqnarray*}
e
\begin{eqnarray*}
\lnot S(\num n) & \lif & T \vdash \lnot !D_S(\num n) \\
& \lif & T \not\vdash !D_S(\num n) \quad \text{(pois $\Th{T}$ é
  consistente)} \\
& \lif & \lnot R(\Gn{!D_S(u)}, n).
\end{eqnarray*}
Isto é, para todo $y$, $S(y)$ é verdadeira se, e somente se,
$R(\Gn{!D_S(u)}, y)$ é verdadeira. Logo, $R$ é universal, e temos a
contradição que procurávamos.
\end{proof}

Seja a ``aritmética verdadeira'' a teoria
$\Setabs{!A}{ \Sat{\Nat}{!A}}$, isto é, o conjunto das sentenças na
linguagem da aritmética que são verdadeiras na interpretação padrão.

\begin{cor}
A aritmética verdadeira não é decidível.
\end{cor}

%Na Seção~\olref{undefinability:truth:section}, enunciaremos um resultado
%mais forte, devido a Tarski.

\end{document}

